\subsection{\texorpdfstring{Relations between the spaces \(L_{F}^{p}\) ( \(1 \leqslant p < +\infty\) )}{Relations between the spaces L sub F to the p ( 1 <= p < +infinity )}}

For every real number \(\alpha > 0\) , the mapping \(z \mapsto |z|^{\alpha-1} \cdot z\) is defined and continuous on the complement of 0 in F; moreover, since \(\left||z|^{\alpha-1} \cdot z\right| = |z|^{\alpha}\) , this function tends to 0 with z and may therefore be extended by continuity to the point 0 by giving it the value 0 at this point, even if \(\alpha < 1\) .

THEOREM 7. --- Let p and q be two real numbers such that \(1 \leqslant p < +\infty\) , \(1 \leqslant q < +\infty\) . If a function f belongs to \(L_{F}^{p}\) then the function \(|f|^{(p/q)-1} \cdot f\) belongs to \(L_{F}^{q}\) , and conversely.

By hypothesis, there exists a sequence \((\mathbf{f}_{n})\) of continuous functions with compact support such that \(\sum_{n=1}^{\infty}\mathrm{N}_{p}(\mathbf{f}_{n})<+\infty\) and \(\mathbf{f}(x)=\sum_{n=1}^{\infty}\mathbf{f}_{n}(x)\) almost everywhere (No.~4, Th. 3). Set

\[
\mathbf {g} _ {n} = | \mathbf {f} _ {1} + \mathbf {f} _ {2} + \dots + \mathbf {f} _ {n} | ^ {(p / q) - 1} \cdot (\mathbf {f} _ {1} + \mathbf {f} _ {2} + \dots + \mathbf {f} _ {n});
\]

the function \(g_{n}\) is continuous with compact support; on the other hand,

\[
\left| \mathbf {g} _ {n} \right| ^ {q} = \left| \mathbf {f} _ {1} + \mathbf {f} _ {2} + \dots + \mathbf {f} _ {n} \right| ^ {p} \leqslant \left(\sum_ {n = 1} ^ {\infty} \left| \mathbf {f} _ {n} \right|\right) ^ {p} = h ^ {q},
\]

where the numerical function \(h \geqslant 0\) (finite or not) satisfies the inequality

\[
\left(\mathrm{N} _ {q} (h)\right) ^ {q} = \left(\mathrm{N} _ {p} \left(\sum_ {n = 1} ^ {\infty} | \mathbf {f} _ {n} |\right)\right) ^ {p} \leqslant \left(\sum_ {n = 1} ^ {\infty} \mathrm{N} _ {p} (\mathbf {f} _ {n})\right) ^ {p} <   + \infty
\]

by the countable convexity theorem. Moreover, \(\mathbf{g}_n(x)\) tends almost everywhere to \(\mathbf{g}(x) = |\mathbf{f}(x)|^{(p/q)-1} \cdot \mathbf{f}(x)\) , therefore Lebesgue's theorem shows that \(\mathbf{g} \in \mathcal{L}_{\mathrm{F}}^q\) . The converse is immediate, since \(\mathbf{f} = |\mathbf{g}|^{(q/p)-1} \cdot \mathbf{g}\) .

It can be shown that the mapping \(\mathbf{f} \mapsto |\mathbf{f}|^{\frac{p}{q}-1} \cdot \mathbf{f}\) is a homeomorphism of \(\mathcal{L}_{\mathrm{F}}^{p}\) onto \(\mathcal{L}_{\mathrm{F}}^{q}\) (§6, Exer. 10).

COROLLARY 1. --- For a function \(\mathbf{f}\) to belong to \(\mathcal{L}_{\mathrm{F}}^{p}\) , it is necessary and sufficient that the function \(|\mathbf{f}|^{p-1} \cdot \mathbf{f}\) belong to \(\mathcal{L}_{\mathrm{F}}^{1}\) .

COROLLARY 2. --- For a positive numerical function f to belong to \(L^{p}\) , it is necessary and sufficient that \(f^{p}\) belong to \(L^{1}\) .

\pandocbounded{\includegraphics[keepaspectratio]{images/part001_881031c4.jpg}}

Note that if \(f\) is a numerical function of arbitrary sign, such that \(|f|^p\) belongs to \(\mathcal{L}^1\) , \(f\) does not necessarily belong to \(\mathcal{L}^p\) (cf.~§4, Exer. 8).

